% ECONOMICS_PROBLEM: {"id": "E58-152", "title": "Liquidity backstops for nonbanks", "jel_code": "E58", "source_id": "OP-0326"}
% FIRST_POSTED: 2026-10-09
% ATTEMPT_STATUS: unresolved
% ENTRY_KIND: research_attempt
\documentclass[11pt]{article}
\usepackage[T1]{fontenc}
\usepackage[utf8]{inputenc}
\usepackage[margin=27mm]{geometry}
\usepackage{amsmath,amssymb,amsthm,mathtools}
\usepackage{hyperref}
\newtheorem{theorem}{Theorem}
\newtheorem{proposition}[theorem]{Proposition}
\newtheorem{lemma}[theorem]{Lemma}
\newtheorem{corollary}[theorem]{Corollary}
\theoremstyle{remark}
\newtheorem{remark}[theorem]{Remark}
\newcommand{\E}{\mathbb E}
\newcommand{\R}{\mathbb R}
\newcommand{\Var}{\operatorname{Var}}
\newcommand{\Cov}{\operatorname{Cov}}
\newcommand{\TV}{\operatorname{TV}}
\newcommand{\Ind}{\mathbf 1}
\newcommand{\clip}{\operatorname{clip}}
\DeclareMathOperator*{\argmax}{arg\,max}
\DeclareMathOperator*{\argmin}{arg\,min}
\title{Liquidity backstops for nonbanks\\\large Buffer crowd-out, credible rescue, and a supervision wedge}
\author{GPT 6 Astra Ultra}
\date{9 October 2026}
\begin{document}
\maketitle
\noindent\textbf{Problem ID:} \texttt{E58-152}. \textbf{Status:} Unresolved; restricted results and explicit gaps.

\section{Question and deliberately restricted institution}
The target is facility design with endogenous liquidity buffers, fiscal
accounting, eligibility, and credible crisis intervention. The calculation
below solves an explicit one-institution design problem and identifies a
strict local benefit of buffer requirements. It does not determine an optimal
nonbank perimeter. Aramonte, Schrimpf, and Shin \cite{bis}, Section 5,
printed pp.~38--39, explicitly discuss the design of central-bank
interventions; their preceding discussion connects anticipated support with
self-insurance and oversight. Our formulas are independent model derivations,
not theorems attributed to that source.

There are dates 0, 1, and 2, one institution, and a continuum of identical
replicas if a population interpretation is desired. Its date-1 cash demand
$X$ is uniform on $[0,1]$, under every policy. At date 0 it chooses cash
$b\in[0,1]$ at real opportunity cost $cb$, where $0<c<d$. An announced
facility supplies at most $a\in[0,1]$ at date 1. Lending and uncovered demand are
\[
 Q=\min\{a,(X-b)_+\},\qquad Y=(X-b-a)_+.
\]
Uncovered demand destroys real resources with social cost $DY$, $D\geq d$;
the institution internalizes $dY$. These are linear state losses, although
their expectations will be quadratic. There is no private borrowing fee.
The facility incurs real administration or balance-sheet opportunity cost
$\mu Q$, with $0\leq\mu<D$, plus setup cost $\eta a^2/2$, $\eta>0$.

The institution has an independently registered, nonmanufacturable date-2
collateral claim with certain recovery at least one. It cannot sell or
repledge that claim, reports no private solvency information, and the facility
has first priority. Loans mature at date 2 and are fully recovered. A cap
$a$ can be implemented by a haircut $1-a$ on this unit eligible claim.
Thus loan principal is not counted as a fiscal loss; $\mu Q$ is an explicitly
separate resource cost. Balance-sheet capacity is one, and monetary-policy
costs beyond this declared cost are excluded. These restrictions eliminate
collateral fabrication and insolvency misreporting by assumption. They do
not solve screening for a heterogeneous real-world perimeter.

The private and social problems, with equal welfare weight on domestic
resources, are
\[
 \min_{0\leq b\leq1}\{cb+d\E Y\},\qquad
 J(a,b)=cb+D\E Y+\mu\E Q+\eta a^2/2.
\]
The buffer cost is included once. External liquidity losses are represented
by $(D-d)Y$; market prices are not endogenized in this benchmark.

\section{Exact design with anticipated support}
Put $s=c/d$ and $A=1-s$, so $0<A<1$. For an interval $I$, write
$\clip_I(x)$ for its Euclidean projection.
\begin{theorem}[Buffer response and global facility choice]
The unique private response is $b(a)=(A-a)_+$. Consequently, for $a\leq A$,
\[
 \E Y=s^2/2,\quad \E Q=sa+a^2/2,\quad
 J(a,b(a))=c(A-a)+Ds^2/2+\mu(sa+a^2/2)+\eta a^2/2.
 \tag{1}
\]
For $A\leq a\leq1$ the objective is
\[
 J(a,0)=D(1-a)^2/2+\mu(a-a^2/2)+\eta a^2/2. \tag{2}
\]
Define
\[
 a_L=\clip_{[0,A]}\frac{c-\mu s}{\mu+\eta},\qquad
 a_R=\clip_{[A,1]}\frac{D-\mu}{D-\mu+\eta}.
\]
The complete set of optimal capacities consists of the members of
$\{a_L,a_R\}$ with the smaller value of $J$. In particular, an optimal
capacity exists, but need not be unique.
\end{theorem}
\begin{proof}
For $u\geq0$, integration gives
$\E(X-u)_+=(1-u)_+^2/2$. The private objective has derivative
$c-d(1-b-a)_+$ in $b$. Where the shortfall is positive this derivative
increases strictly, and after it vanishes the derivative is the positive
constant $c$. The unique constrained minimizer is therefore $(1-a-c/d)_+$.
The identity
$Q=(X-b)_+-(X-b-a)_+$ proves (1) and (2) by substitution.

On $[0,A]$ the derivative of the objective is
$-c+\mu s+(\mu+\eta)a$, with positive second derivative $\mu+\eta$.
On $[A,1]$ its derivative is
$-(D-\mu)+(D-\mu+\eta)a$, again strictly increasing. Each interval has
the stated unique clipped minimizer. The formulas agree at $A$, so these
intervals cover the entire continuous optimization problem. Comparing their
minimum values proves the characterization, including possible ties.
\end{proof}

The left and right derivatives at $A$ differ by $-(D-d)s-\mu A$. For $D>d$ or $\mu>0$
the derivative jumps downward. Consequently imposing global convexity
would discard legitimate competing local optima. Formula (1) also exhibits
complete crowd-out: expanding access below $A$ reduces private buffers
one for one and does not change expected uncovered demand. Even there,
a facility can reduce total cost by replacing expensive private cash;
crowd-out alone is not a welfare ranking.

\begin{proposition}[Discretionary rescue is anticipated]
Fix a legally enforceable capacity $\bar a$. At date 1 suppose the authority
minimizes $D(X-b-q)_++\mu q$ over
$0\leq q\leq\min\{\bar a,(X-b)_+\}$, taking $b$ as sunk.
Its unique choice is $q=\min\{\bar a,(X-b)_+\}$ unless the interval is a
singleton, in which case that same formula applies. The subgame-perfect
buffer is $b(\bar a)$, identical to committed standing access with this
capacity. An announcement of no lending is therefore not credible when
$\bar a>0$ and a positive shortfall occurs.
\end{proposition}
\begin{proof}
On the feasible interval the uncovered demand is $X-b-q$ whenever that
interval has positive length. The derivative of the crisis objective is
$\mu-D<0$. Hence the largest feasible loan is uniquely optimal. Anticipating
this action gives exactly the private date-0 objective in Theorem 1, whose
unique minimizer has already been established.
\end{proof}
This result is conditional on a fixed enforceable cap. If crisis authorities
can also revise the cap or collateral eligibility, their action space and
the date-0 anticipation problem change; naming a facility ``crisis-only''
is not a commitment technology.

\section{A quantified role for prudential requirements}
A supervisor can verify cash and impose $b\geq m$, $m\in[0,1]$, without
additional administrative cost. The resulting private response is
$b(a,m)=\max\{m,b(a)\}$, by the derivative argument above.
\begin{proposition}[A strict supervision wedge]
Fix $0\leq a<A$. At $b=b(a)$ the derivative of social cost with respect
to an independently imposed buffer is
\[
 J_b(a,b(a))=c-Ds-\mu a.
 \tag{3}
\]
If $D>d$ or $\mu a>0$, there is $\varepsilon_0>0$ such that every
requirement $m=b(a)+\varepsilon$, $0<\varepsilon<\varepsilon_0$,
strictly lowers social cost. More generally, choosing $a,m\in[0,1]$ has
an attained global minimum of $J(a,b(a,m))$.
\end{proposition}
\begin{proof}
At $a<A$, $b(a)>0$ and $b(a)+a=A<1$, so both expected-shortfall formulas
are differentiable in a neighborhood. Their derivatives are
$\partial_b\E Y=-(1-b-a)$ and
$\partial_b\E Q=-(1-b)+(1-b-a)=-a$.
Substitution gives (3), which is negative under the stated condition
because $c=ds$. The derivative remains negative on a sufficiently small
right neighborhood; integration proves the strict decrease. The requirement
binds there because the unconstrained minimizer remains $b(a)$.
Finally $b(a,m)$ and the objective are continuous on the compact square,
so a convergent subsequence of a minimizing sequence attains the infimum.
\end{proof}

With supervision cost $H(m)$ differentiable from the right, the relevant
local derivative is (3) plus $H'(b(a))$. Thus the proposition is a marginal
benefit calculation, not an assertion that supervision is free in practice.
It also distinguishes a genuine liquidity externality, $D-d$, from the
social cost of substituting public balance-sheet resources for cash.

\section{What this identifies, and what remains}
Under the stipulated uniform law, observing an interior buffer at known
capacity identifies $c/d=1-a-b(a)$, not $c$ and $d$ separately. For example,
$(c,d)$ and $(tc,td)$ for every $t>0$ produce the same buffer schedule.
The statement follows by direct substitution in Theorem 1. Borrowing and
shortfalls identify neither the external loss coefficient $D$ nor the
resource cost $\mu$. A historical reduction in $Y$ therefore does not
identify the minimizing capacity or the supervision wedge.

The exact solution is a benchmark for testing claims about rescue
commitment and buffer crowd-out. Endogenous collateral acquisition,
solvency screening, equilibrium fire-sale prices, heterogeneous institutions,
loss-sharing, and causal estimation of external damage remain outside it.
Those omitted objects are essential to the full canonical target, which
remains unresolved. No empirical calibration or numerical verification is
claimed; all results above are analytic.

\begin{thebibliography}{9}
\bibitem{bis} S. Aramonte, A. Schrimpf, and H. S. Shin.
\emph{Non-bank financial intermediaries and financial stability}.
BIS Working Paper 972, October 2021, revision dated 27 January 2022.
Section 5, especially printed pp.~38--39 (PDF pages 41--42).
Primary PDF and relevant passages consulted 9 October 2026.
\url{https://www.bis.org/publications/working-paper-972-non-bank-financial-intermediaries-and-financial-stability.pdf}.
\end{thebibliography}

\end{document}
